\chapter{ROS Messages and Services Reference}
\label{app:interfaces}

This appendix documents every message and service defined by
\file{mfja_rail_interfaces}. Table~\ref{tab:interface-source-map} connects each
interface family to its definition and primary runtime implementation. Rebuild
the interface package and all consumers after changing any definition.

\section{Interface Ownership and Source Paths}

The definition column in Table~\ref{tab:interface-source-map} governs the wire
layout. The implementation column identifies the primary runtime source for
the semantics described in the sections that follow. Interface registration
and installation are maintained in
\file{mfja_rail_interfaces/CMakeLists.txt}.

\scriptsize
\begin{longtable}{L{0.18\textwidth} L{0.40\textwidth} L{0.34\textwidth}}
\caption{ROS interface ownership and implementation paths}
\label{tab:interface-source-map}\\
\toprule
\tableheader Interface family & Definition path(s) & Runtime implementation path(s) \\
\midrule
\endfirsthead
\toprule
\tableheader Interface family & Definition path(s) & Runtime implementation path(s) \\
\midrule
\endhead
Shared named state & \file{mfja_rail_interfaces/msg/NamedState.msg} & \file{mfja_robot_control_config/scripts/room_315_kinematic_shuttle_node.py} \\ \rowseparator
Shuttle command and state & \file{mfja_rail_interfaces/msg/ShuttleCommand.msg}; \file{mfja_rail_interfaces/msg/ShuttleState.msg} & \file{mfja_robot_control_config/scripts/room_315_kinematic_shuttle_node.py} \\ \rowseparator
Switch command and state & \file{mfja_rail_interfaces/msg/SwitchCommand.msg}; \file{mfja_rail_interfaces/msg/SwitchState.msg}; \file{mfja_rail_interfaces/msg/NamedState.msg} & \file{mfja_robot_control_config/scripts/room_315_kinematic_shuttle_node.py} \\ \rowseparator
Stopper command and state & \file{mfja_rail_interfaces/msg/StopperCommand.msg}; \file{mfja_rail_interfaces/msg/StopperState.msg}; \file{mfja_rail_interfaces/msg/NamedState.msg} & \file{mfja_robot_control_config/scripts/room_315_kinematic_shuttle_node.py} \\ \rowseparator
Sensor reading and feedback & \file{mfja_rail_interfaces/msg/SensorReading.msg}; \file{mfja_rail_interfaces/msg/SensorFeedback.msg} & \file{mfja_robot_control_config/scripts/room_315_kinematic_shuttle_node.py}; \file{mfja_robot_control_config/scripts/room_315_task_execution.py} \\ \rowseparator
Visual shuttle state & \file{mfja_rail_interfaces/msg/VisualShuttleState.msg} & \file{mfja_robot_control_config/scripts/room_315_visual_state_inference_node.py}; \file{mfja_robot_control_config/scripts/room_315_task_execution.py} \\ \rowseparator
Visual observation & \file{mfja_rail_interfaces/msg/VisualStateObservation.msg}; \file{mfja_rail_interfaces/msg/VisualShuttleState.msg} & \file{mfja_robot_control_config/scripts/room_315_visual_state_inference_node.py}; \file{mfja_robot_control_config/scripts/room_315_task_execution.py} \\ \rowseparator
Add-shuttle service & \file{mfja_rail_interfaces/srv/AddShuttle.srv} & \file{mfja_robot_control_config/scripts/room_315_kinematic_shuttle_node.py} \\
\bottomrule
\end{longtable}
\normalsize

\section{Shared Named State}

\subsection{NamedState.msg}

\begin{lstlisting}
string name
string state
\end{lstlisting}

Switch command input accepts \code{E}/\code{EXTERIOR} or
\code{I}/\code{INTERIOR}; canonical rail state is \code{E} or \code{I}.
Stopper input accepts canonical \code{0} (open/pass) or \code{1}
(closed/stop), with documented textual aliases normalized by the node. Published
state is canonical. Selectors may name an individual configured device or the
supported group \code{ALL}; switch groups also include the runtime's documented
side/group selectors.

\section{Shuttle Command and State}

\subsection{ShuttleCommand.msg}

\begin{lstlisting}
std_msgs/Header header
string name
string command
string start_slot
float64 speed
string target_slot
\end{lstlisting}

\begin{longtable}{L{0.22\textwidth} L{0.68\textwidth}}
\caption{Shuttle command field contract}
\label{tab:shuttle-command-fields}\\
\toprule
\tableheader Field & Contract \\
\midrule
\endfirsthead
\toprule
\tableheader Field & Contract \\
\midrule
\endhead
\code{header} & Command stamp/frame metadata; the rail validates payload fields rather than using the header as completion evidence \\ \rowseparator
\code{name} & Exact Gazebo entity name, empty only for the supported single-shuttle convenience, or \code{ALL} \\ \rowseparator
\code{command} & Canonical operator values ON, OFF, RESET, REMOVE; aliases normalize internally but should not be used in durable clients \\ \rowseparator
\code{start_slot} & Lifecycle/start metadata retained for compatibility; add uses the service request; reset uses configured origin \\ \rowseparator
\code{speed} & Positive configured speed applied only to ON; zero means preserve current ON setting for applicable commands \\ \rowseparator
\code{target_slot} & Optional slot 1--4, legal only with ON; empty means continuous motion \\
\bottomrule
\end{longtable}

\subsection{ShuttleState.msg}

\begin{lstlisting}
std_msgs/Header header
string name
string mode
string current_segment
float64 s
float64 x
float64 y
float64 z
float64 yaw
float64 speed
string reached_target_slot
\end{lstlisting}

Legal modes are \code{MOVING}, \code{WAITING}, \code{DISABLED}, and
\code{FALLING}. \code{s} is arc length in metres on
\code{current_segment}; \code{x,y,z,yaw} are transformed Gazebo pose values.
\code{speed} is retained configured travel speed, not instantaneous velocity.
\code{reached_target_slot} is nonempty only after the controller clamps an
explicit target and disables the shuttle. With no managed shuttle, the node
publishes a blank-name heartbeat to mean known-empty.

\section{Switch and Stopper Arrays}

\subsection{SwitchCommand.msg and SwitchState.msg}

\begin{lstlisting}
std_msgs/Header header
NamedState[] switches
\end{lstlisting}

Command and state share the same structure but not the same meaning: command is
a requested target; state is delayed actual rail authority.

\subsection{StopperCommand.msg and StopperState.msg}

\begin{lstlisting}
std_msgs/Header header
NamedState[] stoppers
\end{lstlisting}

Canonical state \code{0} permits passage and \code{1} stops/blocks. Published
state changes only when the scheduled delay expires in simulation time.

\section{Sensor Feedback}

\subsection{SensorReading.msg}

\begin{lstlisting}
string name
string sensor_type
uint8 active
string shuttle_name
string segment
float64 s
float64 s_ratio
\end{lstlisting}

\code{sensor_type} is currently always \code{sensor}; purpose is encoded by
the configured name/documentation. \code{active} is 0 clear or 1 occupied.
When active, \code{shuttle_name} identifies the occupant. Segment and position
are privileged simulator metadata. The task gateway deliberately ignores those
three position fields and uses only fresh active/name/identity for exact DZI
slot grounding.

\subsection{SensorFeedback.msg}

\begin{lstlisting}
std_msgs/Header header
SensorReading[] readings
\end{lstlisting}

Each side publishes its complete configured/derived reading array at the
configured sensor rate.

\section{Visual Shuttle State}

\subsection{VisualShuttleState.msg}

\begin{lstlisting}
string identity
string presence_state
bool visual_facts_valid
string side
string block
float32[4] bbox_xywh
float64 s_m
float64 s_ratio
float64 segment_length_m
string loaded_state
float64 segment_confidence
float64 loaded_confidence
\end{lstlisting}

Canonical identities are \code{L1--L4,R1--R4}. Presence is
\code{present}, \code{absent}, or \code{unknown}. Valid learned side is
\code{left}/\code{right}; block is one of the 14 public segment classes;
loaded state is \code{empty}/\code{loaded}. The bounding box order is
$(x,y,w,h)$ according to the versioned visual contract. Consumers must honor
\code{visual_facts_valid}; absent/unknown entries do not carry usable learned
facts merely because numeric fields exist.

\section{Visual Observation}

\subsection{VisualStateObservation.msg}

\begin{lstlisting}
std_msgs/Header header
builtin_interfaces/Time left_image_stamp
builtin_interfaces/Time right_image_stamp
string schema_version
string checkpoint_sha256
string stage
bool accepted
bool stabilized
bool stale
bool model_ready
bool input_ready
bool presence_ready
bool state_fusion_ready
string[] validation_reasons
string[] clamped_fields
float64 preprocessing_latency_ms
float64 inference_latency_ms
float64 validation_latency_ms
float64 cycle_latency_ms
uint64 accepted_frame_count
uint64 rejected_frame_count
uint64 stale_frame_count
VisualShuttleState[] shuttles
\end{lstlisting}

\code{stage} distinguishes raw/validation/accepted runtime output. An execution
consumer verifies timestamps, schema, checkpoint, readiness, acceptance,
staleness, and task-required shuttle facts together. \code{stabilized} is a
diagnostic that says the optional temporal filter was applied; it is normally
false on valid accepted observations when that filter is disabled and is not a
task-gateway required-true field. Counters are process-lifetime diagnostics and
are not task sequence numbers.

\section{Add Shuttle Service}

\subsection{AddShuttle.srv}

\begin{lstlisting}
string name
string start_slot
float64 speed
bool start_enabled
---
bool success
string message
string name
\end{lstlisting}

An empty name selects the next unused entity and an empty slot selects the next
available configured slot. A speed that compares greater than zero is used;
otherwise the configured default is substituted, with no separate finiteness
validation. Duplicate entities, invalid slots, and---when configured---occupied
start slots are rejected. A successful response means the logical runtime
accepted the shuttle. Gazebo spawn remains asynchronous, so visual existence
requires separate verification. Removal, reset, ON, and OFF are ShuttleCommand
actions, not services.

\section{Compatibility Rules}

Adding a field is still a ROS interface change and requires dependent
rebuilds. Renaming, removing, retyping, reordering fixed arrays, or changing
units/legal strings is a contract migration. Record a schema/version strategy
for persisted JSON/datasets separately; ROS message compatibility alone does
not migrate recorded episodes or manifests.

Interface registration and dependency declarations are maintained in
\file{mfja_rail_interfaces/CMakeLists.txt}; field definitions are maintained at
the paths in Table~\ref{tab:interface-source-map}.
